\usepackage[tmargin=2cm,rmargin=1in,lmargin=1in,margin=0.85in,bmargin=2cm,footskip=.2in]{geometry}
\usepackage{amsmath,amsfonts,amsthm,amssymb,mathtools}
\usepackage[varbb]{newpxmath}
\usepackage{microtype}
\usepackage{xfrac}
\usepackage{xparse}
\usepackage[makeroom]{cancel}
\usepackage{bookmark}
\usepackage{enumitem}
\usepackage{caption}
\usepackage{hyperref,theoremref}
\hypersetup{
	pdftitle={assignment},
	colorlinks=true, linkcolor=doc!90,
	bookmarksnumbered=true,
	bookmarksopen=true
}
\captionsetup{
	font=small,
	labelfont={bf,sf},
	textfont=sf,
	labelsep=period
}
\usepackage[nameinlink,noabbrev]{cleveref}
\usepackage[most,many,breakable]{tcolorbox}
\usepackage{xcolor}
\usepackage{listings}
\usepackage{varwidth}
\usepackage{etoolbox}
%\usepackage{authblk}
\usepackage{nameref}
\usepackage{multicol,array}
\usepackage[ruled,vlined,linesnumbered]{algorithm2e}
\usepackage{comment} % enables the use of multi-line comments (\ifx \fi) 
\usepackage{import}
\usepackage{xifthen}
\usepackage{pdfpages}
\usepackage{transparent}
\usepackage{chngcntr}
\usepackage{tikz}
\usepackage{titletoc}

\newcommand{\incfig}[1]{%
    \def\svgwidth{\columnwidth}
    \import{./figures/}{#1.pdf_tex}
}

\usepackage{tikzsymbols}
\tikzset{
	symbol/.style={
			draw=none,
			every to/.append style={
					edge node={node [sloped, allow upside down, auto=false]{$#1$}}}
		}
}
\tikzstyle{c} = [circle,fill=black,scale=0.5]
\tikzstyle{b} = [draw, thick, black, -]
\tikzset{
    vertex/.style={
        circle,
        draw,
        minimum size=6mm,
        inner sep=0pt
    }
}

\renewcommand\qedsymbol{$\Laughey$}

% \usepackage{xparse}
% \usepackage{pgffor}
% \usepackage{emoji}
%
% \NewDocumentCommand{\memoji}{m o}{%
%   \IfNoValueTF{#2}
%     {\ifmmode \text{\emoji{#1}} \else \emoji{#1}\fi} % Single emoji case
%     {\foreach \index in {1,...,#2}{\ifmmode \text{\emoji{#1}} \else \emoji{#1}\fi}} % Repeated emoji case
% }

%%%%%%%%%%%%%%%%%%%%%%%%%%%%%%
% SELF MADE COLORS
%%%%%%%%%%%%%%%%%%%%%%%%%%%%%%

\definecolor{doc}{RGB}{0,60,110}
\definecolor{myg}{RGB}{56, 140, 70}
\definecolor{myb}{RGB}{45, 111, 177}
\definecolor{myr}{RGB}{199, 68, 64}
\definecolor{mytheorembg}{HTML}{F2F2F9}
\definecolor{mytheoremfr}{HTML}{00007B}
\definecolor{mylemmabg}{HTML}{FFFAF8}
\definecolor{mylemmafr}{HTML}{983b0f}
\definecolor{mypropbg}{HTML}{f2fbfc}
\definecolor{mypropfr}{HTML}{191971}
\definecolor{myexamplebg}{HTML}{F2FBF8}
\definecolor{myexamplefr}{HTML}{88D6D1}
\definecolor{myexampleti}{HTML}{2A7F7F}
\definecolor{mydefinitbg}{HTML}{FFF5F5}
\definecolor{mydefinitfr}{HTML}{8A1C1C}
\definecolor{myremarkbg}{HTML}{FAF9F9}
\definecolor{myremarkfr}{HTML}{000000}
\definecolor{notesgreen}{RGB}{0,162,0}
\definecolor{myp}{RGB}{197, 92, 212}
\definecolor{mygr}{HTML}{2C3338}
\definecolor{myred}{RGB}{127,0,0}
\definecolor{myyellow}{RGB}{169,121,69}
\definecolor{myexercisebg}{HTML}{F2FBF8}
\definecolor{myexercisefg}{HTML}{88D6D1}
\definecolor{codebg}{HTML}{F7F8FA}
\definecolor{codegutter}{HTML}{EDEFF3}
\definecolor{codeframe}{HTML}{D8DEE8}
\definecolor{codeaccent}{HTML}{3867D6}
\definecolor{codekeyword}{HTML}{7C3AED}
\definecolor{codestring}{HTML}{0F766E}
\definecolor{codecomment}{HTML}{64748B}
\definecolor{codenumber}{HTML}{94A3B8}
\definecolor{codeidentifier}{HTML}{1F2937}
\definecolor{algorule}{HTML}{CBD5E1}
\definecolor{algoaccent}{HTML}{0F766E}
\definecolor{algokeyword}{HTML}{0F172A}
\definecolor{algocomment}{HTML}{64748B}
\definecolor{algonumber}{HTML}{64748B}

%%%%%%%%%%%%%%%%%%%%%%%%%%%%
% ALGORITHM SETUP
%%%%%%%%%%%%%%%%%%%%%%%%%%%%

\newcommand{\cmtalgokw}[1]{\textcolor{algokeyword}{\bfseries\sffamily #1}}
\newcommand{\cmtalgoarg}[1]{\emph{#1}}
\newcommand{\cmtalgofunc}[1]{\textsc{#1}}
\newcommand{\cmtalgodata}[1]{\textsf{#1}}
\newcommand{\cmtalgonumber}[1]{\textcolor{algonumber}{#1}}
\newcommand{\mycommfont}[1]{\footnotesize\sffamily\itshape\textcolor{algocomment}{#1}}

\SetAlFnt{\small}
\SetKwSty{cmtalgokw}
\SetArgSty{cmtalgoarg}
\SetFuncSty{cmtalgofunc}
\SetDataSty{cmtalgodata}
\SetCommentSty{mycommfont}
\SetNoFillComment
\SetNlSty{cmtalgonumber}{\scriptsize\sffamily}{}
\SetAlgoNlRelativeSize{-1}
\SetAlgoInsideSkip{smallskip}
\SetAlgoSkip{medskip}
\SetAlgoHangIndent{1.2em}
\SetAlCapFnt{\small\sffamily\bfseries}
\SetAlCapNameFnt{\small\sffamily}
\SetAlgoCaptionSeparator{:}
\setlength{\algoheightrule}{0.6pt}
\setlength{\algotitleheightrule}{0.6pt}
\setlength{\interspacetitleruled}{4pt}
\setlength{\interspacealgoruled}{4pt}
\crefname{algocf}{algorithm}{algorithms}
\Crefname{algocf}{Algorithm}{Algorithms}

\makeatletter
\newcommand{\cmtalgohline}[2]{%
	\begingroup\color{algorule}\hrule width#1 height#2 depth0pt\endgroup}
\renewcommand{\algocf@Hlne}{%
	\begingroup\color{algorule}\hrule height 0.4pt depth0pt width .5em\endgroup}
\renewcommand{\algocf@Vline}[1]{%
	\strut\par\nointerlineskip
	\algocf@push{\skiprule}%
	\hbox{{\color{algorule}\vrule}%
		\vtop{\algocf@push{\skiptext}%
			\vtop{\algocf@addskiptotal #1}\Hlne}}\vskip\skiphlne
	\algocf@pop{\skiprule}%
	\nointerlineskip}
\renewcommand{\algocf@Vsline}[1]{%
	\strut\par\nointerlineskip
	\algocf@bblockcode
	\algocf@push{\skiprule}%
	\hbox{{\color{algorule}\vrule}%
		\vtop{\algocf@push{\skiptext}%
			\vtop{\algocf@addskiptotal #1}}}%
	\algocf@pop{\skiprule}%
	\algocf@eblockcode}
\def\@algocf@pre@ruled{%
	\cmtalgohline{\algocf@ruledwidth}{\algoheightrule}\kern\interspacetitleruled}
\def\@algocf@post@ruled{%
	\kern\interspacealgoruled\cmtalgohline{\algocf@ruledwidth}{\algoheightrule}}
\renewcommand{\algocf@caption@ruled}{%
	\box\algocf@capbox\kern\interspacetitleruled
	\cmtalgohline{\algocf@ruledwidth}{\algotitleheightrule}\kern\interspacealgoruled}
\makeatother

%%%%%%%%%%%%%%%%%%%%%%%%%%%%
% TCOLORBOX SETUPS
%%%%%%%%%%%%%%%%%%%%%%%%%%%%

\setlength{\parindent}{1cm}
\tcbuselibrary{theorems,skins,hooks}

\tcbset{
	cmtleftpanel/.style 2 args={
		enhanced,
		breakable,
		colback = #1,
		frame hidden,
		boxrule = 0sp,
		borderline west = {2pt}{0pt}{#2},
		sharp corners,
		description font = \mdseries,
		separator sign none,
	},
	cmtleftthm/.style 2 args={
		cmtleftpanel={#1}{#2},
		detach title,
		before upper = \tcbtitle\par\smallskip,
		coltitle = #2,
		fonttitle = \bfseries\sffamily,
		segmentation style={solid, #2},
	},
	cmtexamplethm/.style={
		colback = myexamplebg,
		breakable,
		colframe = myexamplefr,
		coltitle = myexampleti,
		boxrule = 1pt,
		sharp corners,
		detach title,
		before upper=\tcbtitle\par\smallskip,
		fonttitle = \bfseries,
		description font = \mdseries,
		separator sign none,
		description delimiters parenthesis,
	},
	cmtdefinitionthm/.style={
		enhanced,
		before skip=2mm,
		after skip=2mm,
		colback=mydefinitbg,
		colframe=mydefinitfr,
		boxrule=0.5mm,
		attach boxed title to top left={xshift=1cm,yshift*=1mm-\tcboxedtitleheight},
		varwidth boxed title*=-3cm,
		colbacktitle=mydefinitfr,
		coltitle=white,
		boxed title style={
			frame code={
				\path[fill=tcbcolback]
				([yshift=-1mm,xshift=-1mm]frame.north west)
				arc[start angle=0,end angle=180,radius=1mm]
				([yshift=-1mm,xshift=1mm]frame.north east)
				arc[start angle=180,end angle=0,radius=1mm];
				\path[left color=tcbcolback!60!black,right color=tcbcolback!60!black,
					middle color=tcbcolback!80!black]
				([xshift=-2mm]frame.north west) -- ([xshift=2mm]frame.north east)
				[rounded corners=1mm]-- ([xshift=1mm,yshift=-1mm]frame.north east)
				-- (frame.south east) -- (frame.south west)
				-- ([xshift=-1mm,yshift=-1mm]frame.north west)
				[sharp corners]-- cycle;
			},
			interior engine=empty,
		},
		fonttitle=\bfseries,
	},
}

\makeatletter
\tcbset{
	cmtribbonbox/.style={
		enhanced,
		breakable,
		colback=white,
		attach boxed title to top left={yshift*=-\tcboxedtitleheight},
		fonttitle=\bfseries,
		boxed title size=title,
		boxed title style={
			sharp corners,
			rounded corners=northwest,
			colback=tcbcolframe,
			boxrule=0pt,
		},
		underlay boxed title={
			\path[fill=tcbcolframe] (title.south west)--(title.south east)
			to[out=0, in=180] ([xshift=5mm]title.east)--
			(title.center-|frame.east)
			[rounded corners=\kvtcb@arc] |-
			(frame.north) -| cycle;
		},
	},
}
\makeatother

%================================
% THEOREM BOX
%================================

\newtcbtheorem[
	number within=section,
	crefname={theorem}{theorems},
	Crefname={Theorem}{Theorems}
]{Theorem}{Theorem}
{%
	cmtleftthm={mytheorembg}{mytheoremfr},
}
{th}

\newtcbtheorem[
	number within=chapter,
	crefname={theorem}{theorems},
	Crefname={Theorem}{Theorems}
]{theorem}{Theorem}
{%
	cmtleftthm={mytheorembg}{mytheoremfr},
}
{th}


\newtcolorbox{Theoremcon}
{%
	cmtleftpanel={mytheorembg}{mytheoremfr}
}

%================================
% Corollery
%================================
\newtcbtheorem[
	number within=section,
	crefname={corollary}{corollaries},
	Crefname={Corollary}{Corollaries}
]{Corollary}{Corollary}
{%
	cmtleftthm={myp!10}{myp!85!black},
}
{th}
\newtcbtheorem[
	number within=chapter,
	crefname={corollary}{corollaries},
	Crefname={Corollary}{Corollaries}
]{corollary}{Corollary}
{%
	cmtleftthm={myp!10}{myp!85!black},
}
{th}


%================================
% LEMMA
%================================

\newtcbtheorem[
	number within=section,
	crefname={lemma}{lemmas},
	Crefname={Lemma}{Lemmas}
]{Lemma}{Lemma}
{%
	cmtleftthm={mylemmabg}{mylemmafr},
}
{th}

\newtcbtheorem[
	number within=chapter,
	crefname={lemma}{lemmas},
	Crefname={Lemma}{Lemmas}
]{lemma}{Lemma}
{%
	cmtleftthm={mylemmabg}{mylemmafr},
}
{th}

%================================
% Exercise
%================================

\newtcbtheorem[
	number within=section,
	crefname={exercise}{exercises},
	Crefname={Exercise}{Exercises}
]{Exercise}{Exercise}
{%
	cmtleftthm={myexercisebg}{myexercisefg},
}
{th}

\newtcbtheorem[
	number within=chapter,
	crefname={exercise}{exercises},
	Crefname={Exercise}{Exercises}
]{exercise}{Exercise}
{%
	cmtleftthm={myexercisebg}{myexercisefg},
}
{th}


%================================
% PROPOSITION
%================================

\newtcbtheorem[
	number within=section,
	crefname={proposition}{propositions},
	Crefname={Proposition}{Propositions}
]{Prop}{Proposition}
{%
	cmtleftthm={mypropbg}{mypropfr},
}
{th}

\newtcbtheorem[
	number within=chapter,
	crefname={proposition}{propositions},
	Crefname={Proposition}{Propositions}
]{prop}{Proposition}
{%
	cmtleftthm={mypropbg}{mypropfr},
}
{th}


%================================
% CLAIM
%================================

\newtcbtheorem[
	number within=section,
	crefname={claim}{claims},
	Crefname={Claim}{Claims}
]{claim}{Claim}
{%
	cmtleftthm={myg!10}{myg!85!black},
	borderline west = {2pt}{0pt}{myg},
}
{th}



%================================
% EXAMPLE BOX
%================================

\newtcbtheorem[
	number within=section,
	crefname={example}{examples},
	Crefname={Example}{Examples}
]{Example}{Example}
{%
	cmtexamplethm,
}
{ex}

\newtcbtheorem[
	number within=chapter,
	crefname={example}{examples},
	Crefname={Example}{Examples}
]{example}{Example}
{%
	cmtexamplethm,
}
{ex}

%================================
% DEFINITION BOX
%================================

\newtcbtheorem[
	number within=section,
	crefname={definition}{definitions},
	Crefname={Definition}{Definitions}
]{Definition}{Definition}{
	cmtdefinitionthm,
	title={#2},#1}{def}
\newtcbtheorem[
	number within=chapter,
	crefname={definition}{definitions},
	Crefname={Definition}{Definitions}
]{definition}{Definition}{
	cmtdefinitionthm,
	title={#2},#1}{def}

%================================
% REMARK BOX
%================================

\newtcbtheorem[
	number within=section,
	crefname={remark}{remarks},
	Crefname={Remark}{Remarks}
]{Remark}{Remark}
{%
	cmtleftthm={myremarkbg}{myremarkfr},
}
{rmk}

\newtcbtheorem[
	number within=chapter,
	crefname={remark}{remarks},
	Crefname={Remark}{Remarks}
]{remark}{Remark}
{%
	cmtleftthm={myremarkbg}{myremarkfr},
}
{rmk}


%================================
% EXERCISE BOX
%================================

\newcounter{questioncounter}
\counterwithin{questioncounter}{chapter}
% \counterwithin{questioncounter}{section}

\newtcbtheorem[
	use counter=questioncounter,
	crefname={question}{questions},
	Crefname={Question}{Questions}
]{question}{Question}{
	cmtribbonbox,
	colframe=myb!80!black,
	title={#2},
	#1
}{def}

%================================
% SOLUTION BOX
%================================

\newtcolorbox{solution}[1][]{
	cmtribbonbox,
	colframe=myg!80!black,
	title=Solution,
	#1
}
\newenvironment{solbox}[1][]{\begin{solution}[#1]}{\end{solution}}

%================================
% Question BOX
%================================

\newtcbtheorem[
	crefname={question}{questions},
	Crefname={Question}{Questions}
]{qstion}{Question}{
	cmtribbonbox,
	colframe=mygr,
	title={#2},
	#1
}{def}

\newtcbtheorem[
	number within=chapter,
	crefname={wrong concept}{wrong concepts},
	Crefname={Wrong Concept}{Wrong Concepts}
]{wconc}{Wrong Concept}{
	breakable,
	enhanced,
	colback=white,
	colframe=myr,
	arc=0pt,
	outer arc=0pt,
	fonttitle=\bfseries\sffamily\large,
	colbacktitle=myr,
	attach boxed title to top left={},
	boxed title style={
			enhanced,
			skin=enhancedfirst jigsaw,
			arc=3pt,
			bottom=0pt,
			interior style={fill=myr}
		},
	#1
}{def}


%================================
% NOTE BOX
%================================

\usetikzlibrary{arrows,calc,shadows.blur}
\tcbuselibrary{skins}
\newtcolorbox{note}[1][]{%
	enhanced jigsaw,
	colback=gray!20!white,%
	colframe=gray!80!black,
	size=small,
	boxrule=1pt,
	title=\textbf{Note:-},
	halign title=flush center,
	coltitle=black,
	breakable,
	drop shadow=black!50!white,
	attach boxed title to top left={xshift=1cm,yshift=-\tcboxedtitleheight/2,yshifttext=-\tcboxedtitleheight/2},
	minipage boxed title=1.5cm,
	boxed title style={%
			colback=white,
			size=fbox,
			boxrule=1pt,
			boxsep=2pt,
			underlay={%
					\coordinate (dotA) at ($(interior.west) + (-0.5pt,0)$);
					\coordinate (dotB) at ($(interior.east) + (0.5pt,0)$);
					\begin{scope}
						\clip (interior.north west) rectangle ([xshift=3ex]interior.east);
						\filldraw [white, blur shadow={shadow opacity=60, shadow yshift=-.75ex}, rounded corners=2pt] (interior.north west) rectangle (interior.south east);
					\end{scope}
					\begin{scope}[gray!80!black]
						\fill (dotA) circle (2pt);
						\fill (dotB) circle (2pt);
					\end{scope}
				},
		},
		#1,
}

%================================
% CODE BLOCKS
%================================

\tcbuselibrary{listings}

\lstdefinestyle{cmtcodebase}{
	basicstyle=\ttfamily\footnotesize,
	keywordstyle=\bfseries\color{codekeyword},
	commentstyle=\itshape\color{codecomment},
	stringstyle=\color{codestring},
	identifierstyle=\color{codeidentifier},
	numberstyle=\scriptsize\color{codenumber},
	showstringspaces=false,
	tabsize=4,
	columns=fullflexible,
	keepspaces=true,
	upquote=true,
	breaklines=true,
	breakatwhitespace=false,
	postbreak=\mbox{\textcolor{codenumber}{$\hookrightarrow$}\space},
}

\lstdefinestyle{cmtcode}{
	style=cmtcodebase,
	numbers=left,
	numbersep=10pt,
	xleftmargin=2.8em,
}

\lstdefinestyle{cmtplaincode}{
	style=cmtcodebase,
	numbers=none,
	xleftmargin=0pt,
}

\lstdefinestyle{cmtconsole}{
	style=cmtplaincode,
	keywordstyle=\color{codeidentifier},
	commentstyle=\color{codecomment},
	stringstyle=\color{codeidentifier},
}

\lstdefinestyle{cmtlatexcode}{
	style=cmtplaincode,
	language={[LaTeX]TeX},
	texcsstyle=*\bfseries\color{codekeyword},
	keywordstyle=\bfseries\color{codekeyword},
	commentstyle=\itshape\color{codecomment},
	moretexcs={
		Cref,cref,metaname,codeinline,codefile,
		dfn,dfnc,rmk,qs,nt,thm,cor,mlemma,mer,mprop,clm,wc,pf,sol,solve,
		KwIn,KwOut,For,If,Return,tcp
	},
	morekeywords={
		Definition,Lemma,theorem,Remark,question,solution,algorithm,
		codeblock,plaincodeblock,latexcodeblock,consoleblock
	},
}

\tcbset{
	cmtcodebox/.style={
		enhanced,
		breakable,
		sharp corners,
		boxrule=0pt,
		frame hidden,
		colback=codebg,
		colframe=codeframe,
		coltitle=codeidentifier,
		fonttitle=\bfseries\sffamily\small,
		colbacktitle=codegutter,
		titlerule=0pt,
		toptitle=4pt,
		bottomtitle=4pt,
		lefttitle=8pt,
		righttitle=8pt,
		left=8pt,
		right=8pt,
		top=6pt,
		bottom=6pt,
		boxsep=0pt,
		before skip=\baselineskip,
		after skip=\baselineskip,
		borderline west={2pt}{0pt}{codeaccent},
		segmentation hidden,
		listing only,
	},
}

\newtcblisting{codeblock}[2][]{%
	cmtcodebox,
	title={#2},
	listing options={style=cmtcode,language=#2},
	#1
}

\newtcblisting{plaincodeblock}[2][]{%
	cmtcodebox,
	title={#2},
	listing options={style=cmtplaincode,language=#2},
	#1
}

\newtcblisting{latexcodeblock}[1][]{%
	cmtcodebox,
	title={LaTeX},
	listing options={style=cmtlatexcode},
	#1
}

\newtcblisting{consoleblock}[1][]{%
	cmtcodebox,
	title={Console},
	listing options={style=cmtconsole},
	borderline west={2pt}{0pt}{mygr},
	#1
}

\newtcbinputlisting{\codefile}[3][]{%
	cmtcodebox,
	title={#3},
	listing file={#2},
	listing options={style=cmtcode,language=#3},
	#1
}

\DeclareTotalTCBox{\codeinline}{v}{
	nobeforeafter,
	tcbox raise base,
	enhanced,
	boxrule=0pt,
	arc=1pt,
	boxsep=0pt,
	left=3pt,
	right=3pt,
	top=1pt,
	bottom=1pt,
	colback=codebg,
	colframe=codebg,
	fontupper=\ttfamily\footnotesize,
}{#1}

%%%%%%%%%%%%%%%%%%%%%%%%%%%%%%
% SELF MADE COMMANDS
%%%%%%%%%%%%%%%%%%%%%%%%%%%%%%

\newcommand{\thm}[2]{\begin{Theorem}{#1}{}#2\end{Theorem}}
\newcommand{\cor}[2]{\begin{Corollary}{#1}{}#2\end{Corollary}}
\newcommand{\mlemma}[2]{\begin{Lemma}{#1}{}#2\end{Lemma}}
\newcommand{\mer}[2]{\begin{Exercise}{#1}{}#2\end{Exercise}}
\newcommand{\mprop}[2]{\begin{Prop}{#1}{}#2\end{Prop}}
\newcommand{\clm}[3]{\begin{claim}{#1}{#2}#3\end{claim}}
\newcommand{\wc}[2]{\begin{wconc}{#1}{}\setlength{\parindent}{1cm}#2\end{wconc}}
\newcommand{\thmcon}[1]{\begin{Theoremcon}{#1}\end{Theoremcon}}
\newcommand{\ex}[2]{\begin{Example}{#1}{}#2\end{Example}}
\newcommand{\dfn}[2]{\begin{Definition}{#1}{}#2\end{Definition}}
\newcommand{\dfnc}[2]{\begin{definition}{#1}{}#2\end{definition}}
\newcommand{\rmk}[2]{\begin{Remark}{#1}{}#2\end{Remark}}
\newcommand{\qs}[2]{\begin{question}{#1}{}#2\end{question}}
\newcommand{\pf}[2]{\begin{myproof}[#1]#2\end{myproof}}
\newcommand{\nt}[1]{\begin{note}#1\end{note}}

\newcommand*\circled[1]{\tikz[baseline=(char.base)]{
		\node[shape=circle,draw,inner sep=1pt] (char) {#1};}}
\newcommand\getcurrentref[1]{%
	\ifnumequal{\value{#1}}{0}
	{??}
	{\the\value{#1}}%
}
\newcommand{\getCurrentSectionNumber}{\getcurrentref{section}}
\newenvironment{myproof}[1][\proofname]{%
	\proof[\bfseries #1: ]%
}{\endproof}

\newcommand{\mclm}[2]{\begin{myclaim}[#1]#2\end{myclaim}}
\newenvironment{myclaim}[1][\claimname]{\proof[\bfseries #1: ]}{}
\newenvironment{iclaim}[1][\claimname]{\bfseries #1\mdseries:}{}
\newcommand{\iclm}[2]{\begin{iclaim}[#1]#2\end{iclaim}}

\newcounter{mylabelcounter}

\makeatletter
\newcommand{\setword}[2]{%
	\phantomsection
	#1\def\@currentlabel{\unexpanded{#1}}\label{#2}%
}
\makeatother

% deliminators
\DeclarePairedDelimiter{\abs}{\lvert}{\rvert}
\DeclarePairedDelimiter{\norm}{\lVert}{\rVert}

\DeclarePairedDelimiter{\ceil}{\lceil}{\rceil}
\DeclarePairedDelimiter{\floor}{\lfloor}{\rfloor}
\DeclarePairedDelimiter{\round}{\lfloor}{\rceil}

\newsavebox\diffdbox
\newcommand{\slantedromand}{{\mathpalette\makesl{d}}}
\newcommand{\makesl}[2]{%
\begingroup
\sbox{\diffdbox}{$\mathsurround=0pt#1\mathrm{#2}$}%
\pdfsave
\pdfsetmatrix{1 0 0.2 1}%
\rlap{\usebox{\diffdbox}}%
\pdfrestore
\hskip\wd\diffdbox
\endgroup
}
\newcommand{\dd}[1][]{\ensuremath{\mathop{}\!\ifstrempty{#1}{%
\slantedromand\@ifnextchar^{\hspace{0.2ex}}{\hspace{0.1ex}}}%
{\slantedromand\hspace{0.2ex}^{#1}}}}
\ProvideDocumentCommand\dv{o m g}{%
  \ensuremath{%
    \IfValueTF{#3}{%
      \IfNoValueTF{#1}{%
        \frac{\dd #2}{\dd #3}%
      }{%
        \frac{\dd^{#1} #2}{\dd #3^{#1}}%
      }%
    }{%
      \IfNoValueTF{#1}{%
        \frac{\dd}{\dd #2}%
      }{%
        \frac{\dd^{#1}}{\dd #2^{#1}}%
      }%
    }%
  }%
}
\providecommand*{\pdv}[3][]{\frac{\partial^{#1}#2}{\partial#3^{#1}}}
%  - others
\DeclareMathOperator{\Lap}{\mathcal{L}}
\DeclareMathOperator{\Var}{Var} % varience
\DeclareMathOperator{\Cov}{Cov} % covarience
\DeclareMathOperator{\E}{E} % expected

% Since the amsthm package isn't loaded

% I prefer the slanted \leq
\let\oldleq\leq % save them in case they're every wanted
\let\oldgeq\geq
\renewcommand{\leq}{\leqslant}
\renewcommand{\geq}{\geqslant}

%%%%%%%%%%%%%%%%%%%%%%%%%%%%%%%%%%%%%%%%%%%
% TABLE OF CONTENTS
%%%%%%%%%%%%%%%%%%%%%%%%%%%%%%%%%%%%%%%%%%%

\contentsmargin{0cm}
\titlecontents{chapter}[3.7pc]
{\addvspace{30pt}%
	\begin{tikzpicture}[remember picture, overlay]%
		\draw[fill=doc!60,draw=doc!60] (-7,-.1) rectangle (-0.9,.5);%
		\pgftext[left,x=-3.7cm,y=0.2cm]{\color{white}\Large\sc\bfseries Chapter\ \thecontentslabel};%
	\end{tikzpicture}\color{doc!60}\large\sc\bfseries}%
{}
{}
{\;\titlerule\;\large\sc\bfseries Page \thecontentspage
	\begin{tikzpicture}[remember picture, overlay]
		\draw[fill=doc!60,draw=doc!60] (2pt,0) rectangle (4,0.1pt);
	\end{tikzpicture}}%
\titlecontents{section}[3.7pc]
{\addvspace{2pt}}
{\contentslabel[\thecontentslabel]{2pc}}
{}
{\hfill\small \thecontentspage}
[]
\titlecontents*{subsection}[3.7pc]
{\addvspace{-1pt}\small}
{}
{}
{\ --- \small\thecontentspage}
[ \textbullet\ ][]

\makeatletter
\renewcommand{\tableofcontents}{%
	\chapter*{%
	  \vspace*{-20\p@}%
	  \begin{tikzpicture}[remember picture, overlay]%
		  \pgftext[right,x=15cm,y=0.2cm]{\color{doc!60}\Huge\sc\bfseries \contentsname};%
		  \draw[fill=doc!60,draw=doc!60] (13,-.75) rectangle (20,1);%
		  \clip (13,-.75) rectangle (20,1);
		  \pgftext[right,x=15cm,y=0.2cm]{\color{white}\Huge\sc\bfseries \contentsname};%
	  \end{tikzpicture}}%
	\@starttoc{toc}}
\makeatother
